\documentclass[10pt]{article}

\usepackage[margin=0.75in]{geometry}
\usepackage{amsmath,amsthm,amssymb,bm}
\usepackage{xcolor}
\usepackage{cancel}
\usepackage{graphicx}
\usepackage{changepage}
\usepackage{circuitikz}
\usepackage{pgfplots}
\usepackage{physics}
\usepackage{tikz}
\usepackage{tikz-3dplot}
\usepackage{hyperref}
\usepackage{siunitx}
\usepackage{fontspec}
\usepackage{relsize}
\usepackage{subfig}
\usepackage{todonotes}
\usepackage{multicol, multirow, booktabs}
\usepackage[breakable]{tcolorbox}
\usepackage[inline]{enumitem}

\theoremstyle{definition}
\newtheorem{problem}{Problem}
\newtheorem{soln}{Solution}

\pgfplotsset{compat=newest}
\usetikzlibrary{lindenmayersystems}
\usetikzlibrary{arrows}
\usetikzlibrary{calc}
\usetikzlibrary{positioning, fit}
\usetikzlibrary{3d, perspective}
\usetikzlibrary{math}
\usetikzlibrary{fadings}
\usetikzlibrary{angles,quotes} 

\definecolor{incolor}{HTML}{303F9F}
\definecolor{outcolor}{HTML}{D84315}
\definecolor{cellborder}{HTML}{CFCFCF}
\definecolor{cellbackground}{HTML}{F7F7F7}
\newcommand{\ui}{\hat{i}}
\newcommand{\uj}{\hat{j}}
\newcommand{\uk}{\hat{k}}
\newcommand{\ux}{\hat{x}}
\newcommand{\uy}{\hat{y}}
\newcommand{\uz}{\hat{z}}
\newcommand{\uv}[1]{\hat{\bv{{#1}}}}
\newcommand{\pr}[1]{{#1}^\prime}
\pgfdeclarelayer{background}  
\pgfsetlayers{background,main}
\AtBeginDocument{\RenewCommandCopy\qty\SI}
\newcommand{\justif}[2]{&{#1}&\text{#2}}

\makeatletter
\newcommand{\boxspacing}{\kern\kvtcb@left@rule\kern\kvtcb@boxsep}
\makeatother
\newcommand{\prompt}[4]{
    \ttfamily\llap{{\color{#2}[#3]:\hspace{3pt}#4}}\vspace{-\baselineskip}
}

\newcommand{\thevenin}[2]{
  \begin{center}
    \begin{circuitikz} \draw
      (0,0) -- (2,0) to[battery1, l_=$V_{Th}\eq#1$] (2,2) 
      to[resistor, l_=$R_{Th}\eq#2$] (0,2)
      ;
      \draw [o-] (-.07,2.079);
      \draw [o-] (-.07,0.079);
    \end{circuitikz}
  \end{center}
}

\newcommand{\norton}[2]{
  \begin{center}
    \begin{circuitikz} \draw
      (0,0) -- (3,0) to[american current source, l_=$I_{N}\eq#1$] (3,2) -- (0,2) (2,0)
      to[resistor, l=$R_{N}\eq#2$] (2,2)
      ;
      \draw [o-] (-.07,2.079);
      \draw [o-] (-.07,0.079);
    \end{circuitikz}
  \end{center}
}

\DeclareMathOperator{\Div}{div}
\DeclareMathOperator{\Curl}{curl}

\DeclareSIUnit\mile{m}

\newlength\stdbls % standard baselineskip
\AtBeginDocument{%
  \setlength\stdbls{\baselineskip}%
  \gdef\stdblsn{\strip@pt\stdbls}%
} 

\newfontface{\Kaufmann}{Kaufmann}
\DeclareTextFontCommand{\kf}{\Kaufmann}
\newcommand{\scriptr}{\fontsize{12pt}{\stdbls}\kf{r}\fontsize{10pt}{\stdbls}}

\newfontface{\KaufmannB}{Kaufmann Bold}
\DeclareTextFontCommand{\kfb}{\KaufmannB}
\newcommand{\bscriptr}{\fontsize{12pt}{\stdbls}\kfb{r}\fontsize{10pt}{\stdbls}}

\newcommand{\highlight}[1]{\colorbox{yellow}{$\displaystyle #1$}}

\newcommand{\ti}[1]{\widetilde{#1}}

\newcommand{\bv}[1]{\bm{\mathrm{#1}}}

\title{Physics 3130H: Assignment III}
\author{Jeremy Favro (0805980) \\ Trent University, Peterborough, ON, Canada}
\date{\today}

\begin{document}
\maketitle

% PROBLEM 1
\begin{problem}
If the field vector is independent of the radial distance within a sphere, i.e. $\frac{\partial \phi}{\partial r}=C$, find
the function describing the density, $\rho=\rho(r)$.
\end{problem}
\begin{soln}
  Since we consider a spherical mass density and we are given that $\bv{g}$ is constant in $r$ we assume that $\bv{g}=g_r\uv{r}$ and $g_r$ is not a function of $r$ as is given in the question.
  Taking the gradient of $\bv{g}$,
  \begin{align*}
    \bv{\nabla}\cdot\bv{g} & =\frac{1}{r^2}\frac{\partial}{\partial r}\left(r^2 g_r\right) \\
                           & =\frac{2g_r}{r}
  \end{align*}
  and we know from the differential form of Gauss' law for gravitation that this is equal to $-4\pi G\rho(r)$
  and hence we obtain that, because we know that the gravitational field points in the $-\uv{r}$ direction and hence $g_r<0$
  $$\rho(r)=\frac{\abs{g_r}}{2\pi Gr}.$$
\end{soln}

% PROBLEM 2
\begin{problem}
Calculate the gravitational field vector due to a homogeneous cylinder at exterior points on the
axis of the cylinder. Perform the calculation
\begin{enumerate}[label=(\alph*)]
  \item by computing the force directly
  \item by computing the potential first.
\end{enumerate}
\end{problem}
\begin{soln}~
  \begin{enumerate}[label=(\alph*)]
    \item We know that $\bv{F}=-G\frac{Mm}{r^2}\uv{r}$
          and in integral form with is
          $$\bv{F}=-Gm\int_\mathcal{V}\frac{\rho(\bv{r}')}{\scriptr^2}\uv{\bscriptr}d^3r'$$
          where this is an integral over the volume of the cylinder and $\scriptr$ is the distance between a point of mass at $\bv{r}'$ and the point of measurement
          $\bv{r}$. Since we are given that the cylinder is homogenous $\rho(\bv{r}')=\rho$ so the integral becomes
          $$\bv{F}=-Gm\rho\int_0^R\int_0^{2\pi}\int_{0}^{\pi}\frac{1}{\scriptr^2}\uv{\bscriptr} d\theta' d\phi' ds'$$
          here $\scriptr=\abs{\bv{r}-\bv{r}'}=\abs{(0,0,z)-(s'\cos\phi',s'\sin\phi',z')}=\sqrt{s'^2+(z-z')^2}$
          \begin{align*}
            \bv{F} & =-Gm\rho\int_0^R\int_0^{2\pi}\int_{-L}^{L}\frac{(-s'\cos\phi',-s'\sin\phi',z-z')}{\left(s'^2+(z-z')^2\right)^{3/2}}s'dz' d\phi' ds'
          \end{align*}
          Note that the integrals containing $\cos\phi'$ and $\sin\phi'$ go to zero as these functions spend as much time positive as they do negative over the
          $\left[0,2\pi\right]$ interval, so the integral comes down to
          $$\bv{F}=-Gm\rho\int_0^R\int_0^{2\pi}\int_{-L}^{L}\frac{z-z'}{\left(s'^2+(z-z')^2\right)^{3/2}}s'dz' d\phi' ds'\uz$$
          Now we make the substitution $(z-z')^2+s'^2=u \implies du=-2(z-z')dz'$ so
          \begin{align*}
            \bv{F} & =G\frac{m\rho}{2}\int_0^R\int_0^{2\pi}\int_{-L}^{L}\frac{s'}{u^{3/2}}du d\phi' ds'\uz      \\
                   & =-Gm\rho\int_0^R\int_0^{2\pi}\eval{\frac{s'}{\sqrt{s'^2+(z-z')^2}}}_{-L}^{L} d\phi' ds'\uz \\
                   & =-2\pi Gm\rho\int_0^R\frac{s'}{\sqrt{s'^2+(z-L)^2}}-\frac{s'}{\sqrt{s'^2+(z+L)^2}} ds'\uz.
          \end{align*}
          Now we can make the substitution $u=s'^2+(z-L)^2\implies du=2s'ds'$ and similar for $z+L$,
          \begin{align*}
            \bv{F} & =-\pi Gm\rho\left[\int_0^R\frac{1}{\sqrt{u}}du-\frac{1}{\sqrt{u}} du\right]\uz \\
                   & =-2\pi Gm\rho\left[\sqrt{s'^2+(z-L)^2}-\sqrt{s'^2+(z+L)^2}\right]_0^R\uz       \\
                   & =-2\pi Gm\rho\left[\sqrt{R^2+(z-L)^2}-\sqrt{R^2+(z+L)^2}
            -\sqrt{0+(z-L)^2}+\sqrt{0+(z+L)^2}\right]\uz                                            \\
                   & =-2\pi Gm\rho\left[\sqrt{R^2+(z-L)^2}-\sqrt{R^2+(z+L)^2}
            -(z-L)+(z+L)\right]\uz\justif{\quad}{$z>L$}                                             \\
                   & =-2\pi Gm\rho\left[\sqrt{R^2+(z-L)^2}-\sqrt{R^2+(z+L)^2}
              +2L\right]\uz
          \end{align*}
          to obtain $\bv{g}$ we just divide by $m$,
          $$\bv{g}=-2\pi G\rho\left[\sqrt{R^2+(z-L)^2}-\sqrt{R^2+(z+L)^2}
              +2L\right]\uz$$
    \item The gravitational potential for a 3D mass density $\rho$ is given by
          $$\phi=-G\int_{\mathcal{V}}\frac{\rho}{\scriptr}d^3r'$$
          where again $\scriptr$ is the vector between a point of mass at $\bv{r}'$ and a test point $\bv{r}$. Since this is the same vector as previous I will reuse it without rederivation,
          \begin{align*}
            \phi & =-G\int_{\mathcal{V}}\frac{\rho}{\sqrt{s'^2+\left(z-z'\right)^2}}s'ds'd\phi'dz'                                    \\
                 & =-G\rho\int_{-L}^{L}\int_{0}^{2\pi}\int_{0}^R\frac{1}{\sqrt{s'^2+\left(z-z'\right)^2}}s'ds'd\phi'dz'               \\
                 & =-2\pi G\rho\int_{-L}^{L}\int_{0}^R\frac{s'}{\sqrt{s'^2+\left(z-z'\right)^2}}ds'dz'                                \\
                 & =-4\pi G\rho\int_{-L}^{L}\left[\sqrt{s'^2+\left(z-z'\right)^2}\right]_0^Rdz'\justif{\quad}{Same process as before} \\
                 & =-4\pi G\rho\int_{-L}^{L}\sqrt{R^2+\left(z-z'\right)^2}-\left(z-z'\right)dz'.
          \end{align*}
          This is a pretty nasty integral, it involves a trig substitution to end up with the integral of $\sec^3$. I will just use the end
          result:
          \begin{align*}
            \phi & =-2\pi G\rho
            \left[
              -\frac{(z-L)}{2}\sqrt{R^2+(z-L)^2}+\frac{R^2}{2}\ln\left[-2(z-L)+2\sqrt{(z-L)^2+R^2}\right]\right.\\
              &\quad\left.+\frac{(z+L)}{2}\sqrt{R^2+(z+L)^2}-\frac{R^2}{2}\ln\left[-2(z+L)+2\sqrt{(z+L)^2+R^2}\right]-2zL
            \right]                                                                            
          \end{align*}
          and we know that 
          \begin{align*}
            \bv{g}&=-\frac{d\phi}{dz}\uz\\
            &=2\pi G\rho\left[-\frac{R^{2} \left(\frac{2 \left(z + L\right)}{\sqrt{\left(z + L\right)^{2} + R^{2}}} - 2\right)}{2 \left(2 \sqrt{\left(z + L\right)^{2} + R^{2}} - 2 \left(z + L\right)\right)} + \frac{R^{2} \left(\frac{2 \left(z - L\right)}{\sqrt{\left(z - L\right)^{2} + R^{2}}} - 2\right)}{2 \left(2 \sqrt{\left(z - L\right)^{2} + R^{2}} - 2 \left(z - L\right)\right)}\right.\\
             &\quad\left.+ \frac{\sqrt{\left(z + L\right)^{2} + R^{2}}}{2} + \frac{\left(z + L\right)^{2}}{2 \sqrt{\left(z + L\right)^{2} + R^{2}}} - \frac{\sqrt{\left(z - L\right)^{2} + R^{2}}}{2} - \frac{\left(z - L\right)^{2}}{2 \sqrt{\left(z - L\right)^{2} + R^{2}}} - 2L\right]\uz\\
             &=-2\pi G\rho\frac{\sqrt{\left(z + L\right)^{2} + R^{2}} \left(2L \sqrt{\left(z - L\right)^{2} + R^{2}} + z^{2} - 2Lz + R^{2} + L^{2}\right) + \left(-z^{2} - 2Lz - R^{2} - L^{2}\right) \sqrt{\left(z - L\right)^{2} + R^{2}}}{\sqrt{\left(z - L\right)^{2} + R^{2}} \sqrt{\left(z + L\right)^{2} + R^{2}}}\uz\\
             &=-2\pi G\rho\frac{\sqrt{\left(z + L\right)^{2} + R^{2}} \left(2L \sqrt{\left(z - L\right)^{2} + R^{2}} + (z-L)^2 + R^{2} \right) -\left((z+L)^2 + R^{2}\right) \sqrt{\left(z - L\right)^{2} + R^{2}}}{\sqrt{\left(z - L\right)^{2} + R^{2}} \sqrt{\left(z + L\right)^{2} + R^{2}}}\uz\\
             &=-2\pi G\rho\left[\frac{\left(2L \sqrt{\left(z - L\right)^{2} + R^{2}} + (z-L)^2 + R^{2} \right)}{\sqrt{\left(z - L\right)^{2} + R^{2}}}-\frac{\left((z+L)^2 + R^{2}\right)}{ \sqrt{\left(z + L\right)^{2} + R^{2}}}\right]\uz\\
             &=-2\pi G\rho\left[\frac{\left((z-L)^2 + R^{2} \right)}{\sqrt{\left(z - L\right)^{2} + R^{2}}}-\sqrt{\left((z+L)^2 + R^{2}\right)}+2L\right]\uz\\
             &=-2\pi G\rho\left[\sqrt{\left(z - L\right)^{2} + R^{2}}-\sqrt{(z+L)^2 + R^{2}}+2L\right]\uz
          \end{align*}
          which is indeed the same result as when we used the force.
  \end{enumerate}
\end{soln}

% PROBLEM 3
\begin{problem}
Calculate the gravitational potential due to a thin rod of length $l$ and mass $M$ at a perpendicular
distance $R$ from the center of the rod.
\end{problem}
\begin{soln}
  We know that the potential due to a mass distribution is given by the integral of the mass-density function over the extent of the mass distribution, weighted by the mass'
  distance to the measurement point, i.e.
  $$\phi=-G\int_{\mathcal{C}}\frac{\lambda(\bv{r}')}{\scriptr}dl'$$
  where $\scriptr$ takes its usual meaning of the vector between a point of mass and a measurement point, $\bscriptr=\bv{r}-\bv{r}'$, $\lambda(\bv{r}')$ represents a linear mass-density, and we are integrating over a
  contour (closed or open) $\mathcal{C}$. Since in this specific case $\mathcal{C}$ is a straight line we orient our axes so this aligns with $\uz$ and hence $dl'=dz'$. Next we are only given mass
  and not mass density, however mass density for a uniform distribution, which we must assume is the case here, is given by the total mass divided by the space through which it is distributed. Here
  then $\lambda=M/l$. Finally since we are interested in $\scriptr$ and hence $\bscriptr$ we note that since the field is rotationally symmetric about the
  axis of the line we can write $\bv{r}$, the vector from the origin to the measurement point at $R$ as $\bv{r}=R\uv{s}$ and $\bv{r}'=z'\uz$ and hence $\bscriptr=R\uv{s}+0\uv{\phi}-z'\uz$. With this we can
  finally do the integral,
  \begin{align*}
    \phi & =-G \lambda\int_{0}^l\frac{1}{\sqrt{R^2+z'^2}}dz'
  \end{align*}
  because we've done similar integrals a few times now I'll just pull out the formula:
  $$\int \frac{1}{\sqrt{a^2+x^2}}dx=\arcsin\left(\frac{x}{a}\right)$$
  and so
  $$\phi=-G \lambda\eval{\arcsin\left(\frac{z'^2}{R^2}\right)}_{0}^{l}=-\frac{G M}{l}\arcsin\left(\frac{l^2}{R^2}\right).$$
\end{soln}

% PROBLEM 4
\begin{problem}
A particle is dropped into a hole drilled straight through the center of Earth. You assume that
Earth has uniform density. Neglect the rotational effects. Show that the period of the oscillation
is about $\qty{84}{\minute}$.
\end{problem}
\begin{soln}
  The strategy here will be to find the force on the particle as a function of its progress along the hole, then from this find the equations of motion for the particle, solve for
  position as a function of time, then find the period from the obtained function. Since we are assuming Earth's density to be uniform
  we can skip the integrations necessary to derive Newton's law of gravitation and instead use it directly,
  $$\bv{F}=-G\frac{Mm}{r^2}\uv{r}.$$
  However the $M$ term is not the entire mass of the Earth, it is instead the mass inside a \emph{Gaussian surface} of the same radius as the position of the particle.
  Say the radius of the earth is $R$ and the particle position radius is $r<R$ provided the particle is dropped right at the opening of the hole and the Earth is a perfect sphere.
  Then the Gaussian surface encloses $$M_{enc}=4\pi \rho\int_{0}^{r}r'^2dr'=\frac{4\pi r^3\rho}{3}$$
  where $\rho$, the mass density of the Earth, is given by $3M/4\pi R^3$ so $$M_{enc}=\frac{4\pi 3M r^3}{4\pi 3R^3}=M\frac{r^3}{R^3}$$
  and hence
  $$\bv{F}=-G\frac{Mm}{r^2}\uv{r}=-G\frac{Mm}{r^2}\frac{r^3}{R^3}\uv{r}=-G\frac{Mm r}{R^3}\uv{r}.$$
  From here we have a standard dynamics problem. The equation of motion is
  \begin{align*}
             & m\ddot{r}=-G\frac{Mm r}{R^3}    \\
    \implies & \ddot{r}=-G\frac{M}{R^3}r       \\
    \implies & r=Ce^{i\sqrt{G\frac{M}{R^3}}t}.
  \end{align*}
  Here as we are dropping the particle from the surface of the Earth, $C=R$ so that $r(t=0)=R$. Now we could find when this function equals $R$ again to determine the period or we could
  note that what we have is just a simple harmonic oscillator with frequency $\sqrt{G\frac{M}{R^3}}\implies T=2\pi/(\sqrt{G\frac{M}{R^3}})\approx\qty{5096}{\second}=\qty{84.93}{\minute}$.
\end{soln}

% PROBLEM 5
\begin{problem}
A uniform solid sphere of mass $M$ and radius $R$ is fixed a distance $h$ above a thin infinite sheet
of aereal mass density $\rho$. With what force does the sphere attract the sheet?
\end{problem}
\begin{soln}
  Firstly let $\sigma=\rho$ because calling a 2D distribution $\rho$ bugs me. Now the force on any object of mass $M$ placed a distance $h$ above an infinite sheet of mass is given by
  $$\bv{F}=-GM\int_{\mathcal{S}}\frac{\sigma(\bv{r}')}{\scriptr^2}\uv{\bscriptr}da'$$
  with $\bscriptr=-x'\ux-y'\uy+z\uz$. Since the sheet is infinite (and hence is highly symmetric), I will parametrize it in cylindrical coordinates to make the integral easier to do. This means that
  $\bscriptr=-s'\cos\phi'\ux-s'\sin\phi'\uy+z\uz$ and the integral becomes
  $$\bv{F}=-GM\sigma\int_{0}^{\infty}\int_0^{2\pi}\frac{-s'\cos\phi'\ux-s'\sin\phi'\uy+z\uz}{\left(s'^2+z^2\right)^{3/2}}s'd\phi'ds'.$$
  The integrals containing trig functions go to zero and we end up with
  $$\bv{F}=-2\pi GM\sigma\int_{0}^{\infty}\frac{zs'}{\left(s'^2+z^2\right)^{3/2}}ds'\uz$$
  wherein we make the substitution $u=s'^2+z'^2\implies du=2s'ds'$,
  \begin{align*}
    \bv{F} & =-\pi GMz\sigma\int_{0}^{\infty}\frac{1}{\left(u\right)^{3/2}}ds'\uz                 \\
           & =2\pi GMz\sigma \eval{\frac{1}{\sqrt{s'^2+z'^2}}}_0^\infty\uz                        \\
           & =2\pi GMz\sigma \left[\frac{1}{\sqrt{\infty+z'^2}}-\frac{1}{\sqrt{0+z'^2}}\right]\uz \\
           & =-2\pi GM\sigma\uz                                                                   \\
  \end{align*}
  Now since the sheet has infinite mass and therefore infinite intertia, it cannot accelerate. Since the sphere is pinned in place, the system must be in equilibrium as nothing can move. If this is the case then the force of the sphere on the sheet
  must be equal but of opposite sign to the force we calculated of the sheet on the sphere, hence
  $$\bv{F}=2\pi GM\sigma\uz.$$
\end{soln}
\end{document}